% =====================================================================================
% REVISIÓN 1  -  Introducción, Planteamiento del problema, Objetivos, Justificación, Alcances y limitaciones
% Escribe tu texto FUERA de los recuadros \begin{guia} ... \end{guia}.
% No cambies los títulos ni borres los comentarios «% audit:req=...» al final de cada \section / \subsection.
% =====================================================================================

\section{Introducción} \label{sec:introduccion} % audit:req=introduccion

\begin{guia}[Guía: Introducción (Revisión 1)]
\textbf{¿Para qué sirve?} Presenta tu trabajo a un lector que no lo conoce y lo invita a seguir leyendo. Es el resumen
``de entrada'' del proyecto: contexto, problema, propuesta y organización del documento.

\textbf{Debe responder:}
\begin{itemize}[nosep,leftmargin=*]
  \item ¿En qué contexto (sector, empresa, área tecnológica) se desarrolla el proyecto?
  \item ¿Cuál es el problema o la necesidad, dicho en pocas líneas? (el detalle va en \emph{Planteamiento del problema}).
  \item ¿Qué se propone hacer y qué se espera lograr?
  \item ¿Por qué es importante? (una idea; el detalle va en \emph{Justificación}).
  \item ¿Cómo está organizado el documento? (una frase por sección o capítulo).
\end{itemize}

\textbf{Extensión mínima:} 150 palabras (1 a 2 páginas), redactadas con tus propias palabras.

\textbf{Errores frecuentes:} copiar definiciones de Internet; escribir solo ``en este capítulo se verá\ldots'' sin decir
nada concreto; dejar la introducción para el final y no actualizarla cuando el proyecto cambia.

\textbf{Cómo se revisa:} claridad del contexto, que el problema y la propuesta se entiendan sin leer el resto, y que la
organización descrita coincida con las secciones reales del documento.
\end{guia}

En las organizaciones que desarrollan soluciones tecnológicas para la industria, la documentación representa una parte importante del conocimiento operativo. Manuales, instructivos, procedimientos, políticas y especificaciones contienen información que debe consultarse para realizar actividades de manera uniforme. Sin embargo, la existencia de los documentos no garantiza que la información pueda localizarse con rapidez cuando los archivos están distribuidos, tienen formatos distintos o emplean una terminología que no coincide con la pregunta del usuario.

En un entorno de innovación industrial, la información técnica puede cambiar durante el ciclo de vida de un proyecto. Por ello, una persona que necesita conocer un procedimiento requiere una respuesta comprensible y también la posibilidad de revisar el documento, página o sección de donde provino. La trazabilidad permite verificar la respuesta, corregir errores y evitar que una recomendación automática se interprete como una instrucción oficial sin respaldo.

La inteligencia artificial (IA) generativa facilita la interacción con la información mediante preguntas en lenguaje cotidiano. No obstante, un modelo consultado sin documentación específica puede responder con conocimiento general o producir afirmaciones que no están sustentadas por los procedimientos internos. La Recuperación Aumentada por Generación (RAG) combina la búsqueda de información con la generación de texto para proporcionar contexto al modelo y reducir esa limitación \cite{lewis2020rag,gao2023rag}.

Este proyecto propone la base de un copiloto agéntico basado en RAG para documentación técnica y procedimientos de planta. La solución integra una interfaz web, una interfaz de programación de aplicaciones (API) propia y un flujo de IA. El backend recibe la pregunta, valida los datos y prepara la comunicación con el modelo externo, mientras que la arquitectura queda preparada para recuperar fragmentos documentales y devolver sus fuentes.

La interfaz se desarrolla con HTML, CSS y JavaScript; el backend se organiza con Python y FastAPI. La comunicación se realiza mediante HTTP y JSON, y FastAPI permite validar solicitudes y producir documentación interactiva de la API \cite{fastapiDocs}. Las claves de acceso se conservan en el servidor para no exponerlas en el navegador.

El proyecto se realiza en el Centro de Innovación Industrial en Inteligencia Artificial (CII.IA), ubicado en Monterrey, Nuevo León, México. La orientación de la organización hacia la aplicación y el escalamiento de IA para la industria sustenta la pertinencia de una herramienta de consulta documental \cite{ciiia2026,ciiia2026b}. El documento presenta el problema, los objetivos, los antecedentes, la justificación, los alcances y las limitaciones de la propuesta.


\subsection{Planteamiento del problema} \label{sec:problema} % audit:req=planteamiento

\begin{guia}[Guía: Planteamiento del problema (Revisión 1)]
\textbf{¿Para qué sirve?} Describe con precisión la situación que quieres resolver y demuestra que es un problema real.

\textbf{Debe responder:}
\begin{itemize}[nosep,leftmargin=*]
  \item ¿Qué situación actual es insatisfactoria o genera una necesidad? Descríbela con hechos, ejemplos o cifras.
  \item ¿Quién la padece (usuarios, área, empresa) y cómo se manifiesta (retrasos, errores, costos, quejas)?
  \item ¿Qué consecuencias tiene no resolverla?
  \item ¿Qué se ha intentado hasta ahora y qué falta (la brecha)? Conecta con los \emph{Antecedentes}.
  \item ¿Cuál es el problema en una sola frase? Escríbelo como enunciado del problema o como pregunta de
        investigación (p. ej. ``¿Cómo reducir X en Y mediante Z?'').
\end{itemize}

\textbf{Extensión mínima:} 120 palabras.

\textbf{Errores frecuentes:} describir la solución en lugar del problema (``el problema es que no hay un sistema web'');
problemas sin evidencia; enunciados tan amplios que no se pueden resolver en una estadía.

\textbf{Cómo se revisa:} que el problema esté delimitado, sea verificable y que los objetivos lo atiendan.
\end{guia}

El problema se presenta cuando una persona necesita consultar una indicación técnica y no cuenta con un mecanismo uniforme para localizarla. La información puede estar dispersa en documentos distintos, contener versiones previas o emplear términos diferentes a los de la pregunta. La búsqueda manual consume tiempo y puede producir interpretaciones diferentes entre usuarios; una búsqueda literal tampoco siempre identifica fragmentos que expresan la misma idea con otras palabras.

La incorporación directa de un modelo de lenguaje tampoco resuelve por sí sola la situación. Aunque el modelo puede redactar una respuesta clara, no existe garantía de que utilice la documentación interna, conozca la versión vigente o señale la página donde obtuvo la información. En un procedimiento de planta, esta falta de trazabilidad representa una limitación importante.

Por lo tanto, el problema de desarrollo se formula de la siguiente manera: \textquestiondown cómo diseñar e implementar la base de un copiloto web que facilite la consulta de documentación técnica mediante IA, se comunique de manera segura con un backend, utilice un flujo preparado para recuperación semántica y devuelva respuestas con fuentes verificables?

El problema comprende cuatro aspectos relacionados. Primero, la interacción requiere una interfaz sencilla para escribir preguntas y leer respuestas. Segundo, la integración necesita una API que valide y coordine la solicitud. Tercero, el conocimiento requiere un corpus con fragmentos, embeddings y metadatos. Finalmente, el control exige reglas para indicar falta de información, manejar errores y evitar que una respuesta generada se presente como un procedimiento oficial sin evidencia.


\subsection{Objetivos} \label{sec:objetivos} % audit:req=objetivos

\begin{guia}[Guía: Objetivo general y objetivos específicos (Revisión 1)]
\textbf{¿Para qué sirve?} Fija qué vas a lograr. Todo lo demás (resultados y conclusiones) se compara contra estos objetivos.

\textbf{Objetivo general} (uno): verbo en infinitivo + qué se hará + con qué o cómo + para qué, y si es posible un criterio
medible. \emph{Ejemplo:} ``Desarrollar un sistema web de notificaciones para familiares de pacientes hospitalizados que
reduzca en 50\% las llamadas de consulta al área de enfermería.''

\textbf{Objetivos específicos} (de 3 a 6): pasos concretos y verificables que, al cumplirse, logran el objetivo general.
Cada uno debe poder demostrarse con un resultado (un módulo, un diagrama, un experimento) que luego aparecerá en
\emph{Resultados} y se retomará en \emph{Conclusiones}. Deben ser SMART: específicos, medibles, alcanzables, relevantes y
con tiempo definido.

\textbf{Debe responder:} ¿qué voy a entregar? ¿cómo sabré que lo logré? ¿cada objetivo específico aporta al general?

\textbf{Extensión mínima:} 40 palabras en total.

\textbf{Errores frecuentes:} verbos no medibles (``conocer'', ``entender'', ``estudiar''); objetivos que son actividades
(``investigar sobre\ldots'') y no logros; más de un objetivo general; objetivos que luego no se atienden en el documento.

\textbf{Cómo se revisa:} coherencia con el problema, medibilidad y que cada objetivo tenga un resultado asociado.
\end{guia}

\subsubsection{Objetivo general} \label{subsec:objetivo_general}
Desarrollar la base funcional de un copiloto agéntico RAG para consultar documentación técnica y procedimientos de planta, integrando una interfaz web, una API con FastAPI y un modelo externo de IA, con una arquitectura preparada para recuperar fragmentos documentales, conservar sus metadatos y presentar fuentes verificables.

\subsubsection{Objetivos específicos} \label{subsec:objetivos_especificos}
\begin{enumerate}
    \item Analizar la necesidad de consulta documental y delimitar el alcance funcional de la primera versión.
    \item Identificar los actores, entradas, salidas y restricciones que intervienen en una consulta al copiloto.
    \item Diseñar una interfaz web que permita escribir preguntas, enviar solicitudes y visualizar respuestas.
    \item Implementar una API con FastAPI para recibir, validar y responder solicitudes de conversación.
    \item Configurar el intercambio de recursos de origen cruzado (CORS) y el intercambio de datos en formato JSON entre el frontend y el backend.
    \item Integrar un proveedor externo de IA mediante variables de entorno protegidas.
    \item Definir el flujo de validación, recuperación, construcción de contexto y generación de respuestas.
    \item Establecer un contrato para fragmentos documentales con identificador, contenido, archivo, página, sección y puntuación.
    \item Aplicar reglas para que el agente indique cuando no cuenta con información suficiente y evite inventar procedimientos.
    \item Preparar pruebas de comunicación entre frontend, backend, recuperador y modelo externo.
\end{enumerate}


\subsection{Justificación} \label{sec:justificacion} % audit:req=justificacion

\begin{guia}[Guía: Justificación (Revisión 1)]
\textbf{¿Para qué sirve?} Explica por qué vale la pena dedicar tiempo y recursos a este proyecto.

\textbf{Debe responder:}
\begin{itemize}[nosep,leftmargin=*]
  \item ¿Quién se beneficia y de qué manera? (empresa, usuarios, la comunidad, tú).
  \item ¿Qué se gana en términos prácticos, económicos, sociales, tecnológicos o científicos? Usa datos cuando existan.
  \item ¿Qué pasaría si no se hace?
  \item ¿Por qué es viable hacerlo ahora y por qué con este enfoque?
  \item ¿Cómo se relaciona con el perfil de egreso de tu programa académico?
\end{itemize}

\textbf{Extensión mínima:} 100 palabras.

\textbf{Errores frecuentes:} repetir el problema con otras palabras; afirmaciones sin sustento (``mejorará mucho la
eficiencia''); olvidar al beneficiario.

\textbf{Cómo se revisa:} que los beneficios sean concretos y consistentes con el problema y los objetivos.
\end{guia}

Un copiloto para documentación técnica puede reducir el tiempo dedicado a localizar información y facilitar el acceso al conocimiento organizacional. En lugar de depender de que el usuario recuerde el nombre de un archivo o revise manualmente varias carpetas, la interfaz permite formular una pregunta de manera natural y recibir una respuesta vinculada con una fuente.

La propuesta es especialmente útil cuando la documentación contiene procedimientos repetitivos, definiciones técnicas o instrucciones que deben consultarse con frecuencia. La recuperación semántica puede identificar fragmentos relacionados aunque la redacción de la pregunta no coincida literalmente con la del documento. La inclusión de metadatos permite revisar el resultado y confirmar la información antes de utilizarla en una actividad crítica.

Desde el punto de vista de desarrollo, el proyecto integra conocimientos de frontend, backend, API, bases de datos, procesamiento de lenguaje natural y evaluación de IA. La separación de componentes facilita que el diseño visual cambie sin alterar el servicio, que el proveedor de IA pueda sustituirse y que la fuente documental se actualice sin reconstruir toda la aplicación.

La propuesta también se justifica por su enfoque de control. No se plantea un chatbot abierto que responda cualquier tema, sino un copiloto limitado a un dominio documental. Cuando el contexto no sea suficiente, el sistema debe comunicar esa condición; una respuesta negativa y verificable es preferible a una respuesta extensa sin respaldo.


\subsection{Alcances y limitaciones} \label{sec:alcances_limitaciones} % audit:req=alcances

\begin{guia}[Guía: Alcances y limitaciones (Revisión 1)]
\textbf{¿Para qué sirve?} Marca las fronteras del proyecto para evitar expectativas irreales.

\textbf{Alcances} -- lo que \emph{sí} incluye el trabajo. Debe responder: ¿qué módulos, funciones, procesos, usuarios o
datos se cubren? ¿en qué plataformas o entornos? ¿qué entregables habrá? ¿qué queda \emph{explícitamente fuera}?

\textbf{Limitaciones} -- lo que restringe el trabajo y que no depende totalmente de ti. Debe responder: ¿qué límites hay
de tiempo (un cuatrimestre), recursos, acceso a datos o a la infraestructura de la empresa, tecnología, permisos?
¿cómo pueden afectar la validez o generalización de los resultados?

\textbf{Extensión mínima:} 60 palabras en total.

\textbf{Errores frecuentes:} prometer más de lo que cabe en la estadía; confundir limitaciones con excusas; no decir qué
queda fuera.

\textbf{Cómo se revisa:} realismo del alcance y honestidad de las limitaciones.
\end{guia}

\subsubsection{Alcances} \label{subsec:alcances}
\begin{itemize}
    \item Analizar la necesidad de consulta documental y definir el flujo general del copiloto.
    \item Mantener una interfaz web para escribir preguntas, enviar solicitudes y consultar respuestas.
    \item Implementar una API con FastAPI para validar solicitudes y devolver respuestas en JSON.
    \item Configurar la comunicación entre frontend y backend mediante CORS.
    \item Integrar un modelo externo de IA mediante variables de entorno.
    \item Definir la arquitectura RAG, la recuperación de contexto, la construcción del prompt y la generación de respuestas.
    \item Especificar los metadatos mínimos de cada fragmento: identificador, contenido, archivo de origen, página, sección y puntuación.
    \item Preparar la integración con documentos técnicos, embeddings, almacenamiento vectorial y búsqueda semántica.
    \item Documentar las decisiones de diseño y las pruebas de comunicación de la primera versión.
\end{itemize}

\subsubsection{Limitaciones} \label{subsec:limitaciones}
\begin{itemize}
    \item La respuesta depende de la disponibilidad, cuota, autenticación y modelo seleccionado en el proveedor externo de IA.
    \item La calidad del resultado depende de que los documentos estén completos, vigentes, correctamente fragmentados y acompañados de metadatos.
    \item La búsqueda semántica no garantiza por sí sola que un fragmento contenga toda la respuesta; se requiere validar la suficiencia del contexto.
    \item La carga de embeddings y el endpoint de búsqueda deben validarse con los datos autorizados y entregados por el Integrante 1.
    \item El sistema no sustituye la revisión de una persona responsable cuando la consulta se relaciona con seguridad, calidad, operación o una decisión crítica.
    \item La primera versión se limita al dominio de documentación técnica y procedimientos de planta; no es un chatbot de propósito general.
    \item El cambio de proveedor, la actualización del corpus y el despliegue en producción requieren configuraciones y pruebas adicionales.
\end{itemize}
